\documentclass{article}
\usepackage[T1]{fontenc}
\usepackage{lmodern}
\usepackage{graphicx}
\usepackage{amsmath,amssymb}
\usepackage[a4paper,margin=2cm]{geometry}
\usepackage[absolute,overlay]{textpos}
\setlength{\TPHorizModule}{1mm}
\setlength{\TPVertModule}{1mm}
\usepackage[indonesian]{babel}
\usepackage{hyperref}
\setlength{\emergencystretch}{2em}
% unit: Halaman 14

\begin{document}

% === Halaman 14 (python/native) ===

Metode Steganalisis

1 . Serangan berbasis visual (visual attacks) 

• Khusus untuk stego-object berupa citra

• Bersifat subjektif, karena melakukan pengamatan  secara kasat mata dengan 

melihat artefak yang mencurigakan di dalam stego-image, lalu 

membandingkannya dengan citra asli (cover image)

• Digunakan pada masa-masa awal riset steganalisis     

• Contoh serangan visual: 


   a. LSB plane attack


   b. Filtered visual attack (Enhanced LSB)  

14

\end{document}
